\section{Data Sources and Acquisition Methodology}

\subsection{Primary and reference sources}

Daily market observations are requested from Yahoo Finance through
\texttt{yfinance} \citep{yfinance,yahoofinance} with
\texttt{auto\_adjust}\allowbreak\texttt{=False} and corporate actions
enabled. SEC EDGAR filings and reference files support registrant identity,
exchange, ticker, and historical-eligibility evidence. The official exchange
calendar is used to separate scheduled non-trading days from unexpected gaps.
No source is assumed complete or error-free.

\subsection{Batch and retry design}

Requests are processed in configurable batches with per-security checkpoint
files, bounded retries, and conservative request spacing. A failed request does
not invalidate its batch. Every attempt is represented in the manifest or
failure register, including timestamps, parameters, provider/package versions,
observation counts, and partial-history status.

\subsection{Raw-data preservation}

Raw observations are written before cleaning. Validation does not overwrite the
raw vendor files. Standardized outputs retain source and acquisition lineage so
that a validated row can be traced to a specific retrieval event.

\InputGenerated{generated/acquisition_results.tex}

\subsection{Acquisition-provenance reconciliation}

The research-facing dataset is derived exclusively from the Yahoo raw
directory. Its successful security identifiers are a subset of the successful
Yahoo manifest entries, and the EODHD public-demo files remain under a separate
path and provider label. The earlier contradiction arose because the abstract
was hard-coded while the body and tables were generated from the refreshed
Yahoo statistics. The abstract now consumes a generated artifact, and a
machine-readable consistency gate compares report macros with manifest and
base counts. The verified identity tests cover \IdentityValid{} of
\YahooSuccess{} checkpoints.

The restricted market-data cross-check uses the EODHD public-demo API
\citep{eodhd}; its small sample is not a replacement for Yahoo.

\IfFileExists{generated/table_download_status.tex}{
  \begin{table}[htbp]
    \centering
    \caption{Acquisition status for the attempted historical universe.}
    \label{tab:download_status}
    \small
    \begin{tabular}{@{}lrrrr@{}}
\toprule
Provider & Requested & Success & Failed & Partial \\
\midrule
Yahoo Finance (primary) & 4,874 & 3,239 & 1,635 & 0 \\
EODHD public demo (secondary) & 40 & 3 & 37 & 0 \\
\bottomrule
\end{tabular}

  \end{table}
}{\PendingResults{generated/table_download_status.tex}}

\FigureGenerated{../figures/download_status.pdf}
  {Observed successes and failures for the full Yahoo universe attempt and the restricted secondary cross-check.}
  {fig:download_status}
